% Recuperación documental para el artículo V; RESULTADO_RECUPERADO.
% Propietario: /Users/ruben/Documents/New project/output/REVISION_EDITORIAL_HMT_MD_20260905/01_FUENTES/INTEGRAL/tree/New project/output/ENSAMBLADO_CAUSAL_RESTAURADO_HMT_MD_20260905/fuente/manuscrito/sections/md/delta_127/04b_completaciones_cuanticas.tex
% Líneas de origen: 1--334.
% REV02: construcción algebraica; la aplicación estadística se reúne en 40.
% Correspondencia editorial: gestion/CONSOLIDACION_PAULI_FOCK_REV02.json.
\section{Construction algébrique des complétions quantiques sur APP}
\label{v:rec:chap:completaciones-cuanticas}
\label{v:sec:completaciones-algebraicas}

APP et TPK fournissent des états orientés, un transport et une parité de feuillet.
Leur réalisation linéaire permet de construire les algèbres symétrique et extérieure,
avec leurs opérateurs de création, d’annihilation et d’occupation. Cette section
établit ces constructions algébriques, le secteur cellulaire des champs et
le couplage réversible de mesure. Le transport des complétions
par raffinement est démontré dans la
section~\ref{v:rec:fock_gibbs:section:01} ; l’application de leurs spectres
d’occupation aux dénombrements et à l’équilibre est développée dans les
sections~\ref{v:rec:chap:ocupacion-estadistica}--\ref{v:rec:ocupacion:section:03}.
Chaque étape conserve les conditions de la représentation qu’elle utilise.

\subsection{Espace linéaire des états associé à une réalisation de particule}
\label{v:rec:completaciones:section:02}

Soit \(S\) un ensemble fini de classes de chemins TPK à la profondeur choisie
et soit
\[
 \mathcal H_1=\ell^2(S)
\]
l’espace de Hilbert de base orthonormale \((e_s)_{s\in S}\). Un choix de
poids positifs \(E_s\) définit l’opérateur diagonal
\[
 He_s=E_se_s.
\]
Le choix de \(S\), de son quotient par les symétries et des poids fait partie
du modèle. Une fois ce choix fixé, les complétions multiparticulaires sont canoniques.

\subsection{Complétions symétrique et extérieure}
\label{v:rec:completaciones:section:03}

On définit
\[
 \mathcal F_+(\mathcal H_1)=
 \bigoplus_{n\geq0}\operatorname{Sym}^n\mathcal H_1,
 \qquad
 \mathcal F_-(\mathcal H_1)=
 \bigoplus_{n\geq0}\bigwedge^n\mathcal H_1.
\]
La première est la complétion bosonique ; la seconde, la fermionique. Dans les deux,
le terme de somme directe de degré zéro est l’espace du vide. Les puissances extérieures
sont munies du produit scalaire pour lequel les produits extérieurs ordonnés de vecteurs
d’une base orthonormale sont orthonormaux ; les puissances symétriques sont utilisées
avec la normalisation usuelle des opérateurs de création et d’annihilation.

\subsection{Graduation de feuillet et algèbre d’échange}
\label{v:rec:completaciones:section:04}

Supposons que le feuillet orienté définit un caractère additif
\[
 p:\operatorname{Mor}(\mathsf P_{\rm TPK})\longrightarrow\mathbb Z/2\mathbb Z,
 \qquad p(\gamma\eta)=p(\gamma)+p(\eta).
\]
Sur l’espace linéaire gradué correspondant, on introduit la symétrie
\[
 \tau(v\otimes w)=(-1)^{p(v)p(w)}w\otimes v.
\]

\begin{theorem}[Complétion statistique déterminée par la parité]
\label{v:rec:completaciones:statement:02}
Le quotient de l’algèbre tensorielle par les relations
\[
 v\otimes w-(-1)^{p(v)p(w)}w\otimes v
\]
est symétrique sur le secteur pair et extérieur sur le secteur impair. Deux états
impairs identiques ont un produit nul ; les états pairs admettent une occupation
arbitraire.
\end{theorem}

\begin{proof}
Si \(p(v)=p(w)=0\), la relation est \(vw=wv\). Si
\(p(v)=p(w)=1\), elle est \(vw=-wv\) ; en prenant \(v=w\) et en travaillant en
caractéristique zéro, on obtient \(v^2=0\). Les secteurs mixtes commutent sans
signe supplémentaire. C’est la propriété universelle de l’algèbre symétrique
graduée.
\end{proof}

\subsection{Opérateurs canoniques et idempotence de l’occupation}
\label{v:pauli:car-construccion}

Le caractère de parité détermine la complétion ; le produit de cette
complétion détermine ensuite l’algèbre d’opérateurs. Dans le
secteur extérieur, la création \(a_s^\dagger\) est le produit extérieur par
\(e_s\), et \(a_s\) est son adjoint.

\begin{theorem}[Principe d’exclusion de Pauli dans la complétion extérieure]
\label{v:rec:completaciones:statement:01}
\label{v:rec:thm:principio-exclusion-pauli}
Pour tout mode \(s\in S\), les opérateurs de création et d’annihilation de la
complétion extérieure satisfont
\begin{equation}
 \{a_s,a_t^\dagger\}=\delta_{st}I,
 \qquad
 \{a_s,a_t\}=\{a_s^\dagger,a_t^\dagger\}=0.
 \label{v:pauli:eq:car}
\end{equation}
En particulier, l’opérateur d’occupation
\(N_s=a_s^\dagger a_s\) est un projecteur :
\begin{equation}
 N_s^2=N_s,
 \qquad \Spec(N_s)\subseteq\{0,1\}.
 \label{v:pauli:eq:proyeccion-ocupacion}
\end{equation}
Les valeurs propres \(0\) et \(1\) correspondent, respectivement, à l’absence
et à la présence du mode \(s\).
\end{theorem}

\begin{proof}
Sur un produit extérieur ordonné, \(a_s\) supprime le facteur \(e_s\), lorsqu’il est
présent, avec le signe de sa position ; s’il est absent, il donne zéro. Composer
cette opération avec le produit extérieur démontre les relations
\eqref{v:pauli:eq:car} : pour \(s=t\), les deux compositions couvrent les cas
complémentaires d’absence et de présence ; pour \(s\ne t\), les signes
d’échange s’annulent. L’antisymétrie implique
\(e_s\wedge e_s=0\) et
\(a_s^2=(a_s^\dagger)^2=0\). Alors
\begin{equation}
 N_s^2=a_s^\dagger a_s a_s^\dagger a_s
 =a_s^\dagger(I-a_s^\dagger a_s)a_s
 =a_s^\dagger a_s=N_s.
 \label{v:pauli:eq:prueba-idempotencia}
\end{equation}
L’opérateur est autoadjoint, et toute valeur propre d’un idempotent satisfait
\(\lambda^2=\lambda\). Son spectre est donc contenu dans
\(\{0,1\}\). Le vide et l’état \(e_s\) réalisent les deux valeurs
\cite{Pauli1925,Dirac1926}.
\end{proof}

Dans le secteur symétrique, les opérateurs bosoniques \(b_s^\dagger,b_s\)
agissent sur la base normalisée des occupations par les facteurs
\(\sqrt{m_s+1}\) et \(\sqrt{m_s}\), respectivement. De leurs compositions,
on obtient
\begin{equation}
 [b_s,b_t^\dagger]=\delta_{st}I,
 \qquad [b_s,b_t]=[b_s^\dagger,b_t^\dagger]=0.
 \label{v:pauli:eq:ccr}
\end{equation}
Ces identités sont comprises sur le domaine algébrique des particules
finies, stable sous les opérateurs considérés.

Le théorème~\ref{v:rec:completaciones:statement:02} transforme une parité compositionnelle TPK en une règle exacte
d’occupation. Pour obtenir le théorème physique de spin--statistique, il faut
construire en outre une théorie locale lorentzienne positive et démontrer que la
parité de feuillet coïncide avec le spin sous les rotations. L’égalité ne se
déduit pas uniquement du fait que les deux objets prennent leurs valeurs modulo deux.

\subsection{Holonomie de feuillet et périodicité de \texorpdfstring{\(2\pi/4\pi\)}{2pi/4pi}}
\label{v:rec:completaciones:section:05}

La parité précédente admet une réalisation géométrique naturelle lorsqu’une
extension survivante possède un retour fermé dont le revêtement d’orientation
est non trivial. Soit \(\rho\) le générateur du retour et soit
\[
 \varepsilon:\langle\rho\rangle\longrightarrow\{-1,+1\}
\]
son caractère d’holonomie, avec \(\varepsilon(\rho)=-1\). Le fibré réel
associé est alors la droite d’orientation d’une bande de Möbius : un
tour ferme la projection géométrique, mais change le feuillet relevé ; deux
tours restaurent l’orientation.

Pour écrire cette propriété angulairement, fixons une représentation
unitaire \(U\) du revêtement des rotations sur l’espace gradué
\(\mathcal H_1=\mathcal H_{\bar0}\oplus\mathcal H_{\bar1}\), compatible avec
la parité de feuillet au sens
\[
 U(2\pi)v=(-1)^{p(v)}v
\]
pour tout vecteur homogène \(v\). La compatibilité est une condition de la
réalisation : elle identifie le générateur du revêtement d’orientation au
caractère \(p\) déjà construit sur les chemins.

\begin{theorem}[Restauration de l’orientation et algèbre d’échange]
\label{v:rec:thm:dos-cuatro-pi-estadistica}
Sous la compatibilité précédente,
\[
 U(2\pi)|_{\mathcal H_{\bar0}}=I,
 \qquad
 U(2\pi)|_{\mathcal H_{\bar1}}=-I,
 \qquad
 U(4\pi)=I.
\]
En outre, le même caractère \(p\) détermine la symétrie graduée
\[
 \tau(v\otimes w)=(-1)^{p(v)p(w)}w\otimes v.
\]
Par conséquent, le secteur d’holonomie paire est complété au moyen de puissances
symétriques et le secteur d’holonomie impaire au moyen de puissances extérieures. Dans
le second, l’exclusion \(v\wedge v=0\) est satisfaite.
\end{theorem}

\begin{proof}
La première et la seconde égalité sont la condition de compatibilité évaluée
dans chaque composante homogène. En itérant un tour,
\[
 U(4\pi)v=U(2\pi)^2v=(-1)^{2p(v)}v=v,
\]
ce qui démontre la restauration. L’affirmation sur l’échange est le
théorème~\ref{v:rec:completaciones:statement:02} : pour deux vecteurs
pairs, le signe est positif, tandis que pour deux impairs, il est négatif.
\end{proof}

Il convient de distinguer trois retours. La clôture angulaire après \(2\pi\) est un
retour de la projection ; la clôture après \(4\pi\) restaure l’orientation
relevée ; aucune n’oblige à elle seule toutes les coordonnées du registre à
revenir à leur valeur initiale. Cette séparation est l’analogue dimensionnel de la
distinction HMT entre retour de phase et retour d’état.

\subsubsection{Les feuillets somme et produit comme supports des \(2\) complétions}
\label{v:rec:completaciones:section:06}

Le tableau APP contient deux évaluations sur un même support. Le feuillet
additif agit par translations et conserve la multiplicité des
continuations ; en ce sens, il accumule. Le feuillet multiplicatif est
permutatif sur les unités et contracte des fibres sur les non-inversibles ; en
ce sens, il plie. Cette incompatibilité est antérieure au choix
statistique et explique pourquoi TPK doit conserver deux généalogies au lieu de
les réduire à un seul mot visible.

L’attribution statistique ne s’obtient toutefois pas en remplaçant
verbalement « somme » par « boson » et « produit » par « fermion ». L’application exacte est
le caractère de parité
\[
 p:\operatorname{Mor}(\mathsf P_{\rm TPK})\longrightarrow\mathbb Z/2\mathbb Z.
\]
Lorsque le transport d’une extension tombe dans le secteur \(p=0\), sa
complétion admet une occupation cumulative et est symétrique. Lorsqu’il tombe dans
\(p=1\), l’inversion de feuillet introduit le signe d’échange et la
complétion est extérieure. La somme et le produit fournissent les deux
géométries locales ; le registre, l’orientation et le caractère \(p\) décident dans
quel secteur un chemin est réalisé.

La périodicité \(4\pi\) et l’exclusion fermionique sont coordonnées par
le même caractère de feuillet. Leurs conséquences pour les dénombrements finis et
les occupations d’équilibre s’obtiennent, avec le Hamiltonien et les
conditions thermiques explicites, dans la
section~\ref{v:rec:chap:ocupacion-estadistica} et dans le
théorème~\ref{v:rec:thm:factorizacion-gran-canonica-rev3}.

\subsection{Le complexe cellulaire du tore APP}
\label{v:rec:completaciones:section:08}

Le graphe de APP est le squelette de dimension un d’une décomposition cellulaire du tore
\(T^2\) à \(9\times9\) sommets. Soit
\[
 C^0\xrightarrow{d_0}C^1\xrightarrow{d_1}C^2
\]
son complexe de cochaînes réelles. Le produit scalaire uniforme étant choisi,
nous définissons la codifférentielle \(\delta=d^*\) et le laplacien
\[
 \Delta_k=d_{k-1}\delta_k+\delta_{k+1}d_k.
\]

\begin{theorem}[Secteur harmonique du tore numérique]
\label{v:rec:completaciones:statement:04}
La dimension de l’espace des formes harmoniques de degré un est deux :
\[
 \dim\ker\Delta_1=b_1(T^2)=2.
\]
On peut choisir une base représentée par les flux uniformes horizontal et
vertical. Les deux sont fermés, cofermés et ne sont pas des gradients de fonctions
globales périodiques.
\end{theorem}

\begin{proof}
Le théorème de Hodge pour les complexes cellulaires finis identifie
\(\ker\Delta_1\) à \(H^1(T^2;\mathbb R)\). La cohomologie du tore est de
rang deux. Directement, les flux constants horizontal et vertical ont
une divergence et un rotationnel nuls ; leurs intégrales sur les deux cycles
fondamentaux sont indépendantes, de sorte qu’ils ne sont pas exacts.
\end{proof}

L’espace des connexions modulo les transformations de jauge est
\[
 C^1/\operatorname{im}d_0.
\]
Dans un complexe fini muni du produit scalaire précédent, chaque orbite
admet un représentant de Coulomb dans \(\ker\delta_1\) ; à l’intérieur de l’orbite,
le choix est unique modulo le stabilisateur \(\ker d_0\) du paramètre de
jauge. Le quotient cohomologique correct est
\[
 \ker d_1/\operatorname{im}d_0
 \simeq \ker\Delta_1 .
\]
Ses deux classes harmoniques constituent un candidat mathématique naturel pour
deux modes globaux libres. Ce ne sont pas encore les deux polarisations locales d’un
photon dans l’espace-temps \(3+1\) : cette identification nécessite une
limite dynamique, une relation de dispersion et un groupe physique de jauge.

\subsection{Loi de Gauss discrète}
\label{v:rec:completaciones:section:09}

Pour une 1-cochaîne orientée \(F\) et un ensemble de sommets \(V\), la
divergence est définie par
\[
 (\operatorname{div}F)(v)=
 \sum_{e:v\to w}F(e)-\sum_{e:w\to v}F(e).
\]
En sommant sur \(V\), toute arête intérieure apparaît une fois avec chaque signe et
s’annule. Par conséquent,
\[
 \boxed{
 \sum_{v\in V}\operatorname{div}F(v)
 =\operatorname{Flux}_{\partial V}(F).}
\]
Cette identité est une loi de Gauss exacte sur le graphe. Interpréter
\(\operatorname{div}F\) comme une charge électrique exige le caractère de charge et la
réalisation dimensionnelle construits en mécanique dimensionnelle.

\subsection{Appareil de mesure réversible}
\label{v:rec:completaciones:section:10}

Soit \(S\) un ensemble d’états et soit \(r:S\to\mathbb Z/9\mathbb Z\) une
lecture. Le couplage à un registre nonadique est défini par
\[
 U_r:S\times\mathbb Z/9\mathbb Z\longrightarrow
 S\times\mathbb Z/9\mathbb Z,
 \qquad
 U_r(s,a)=(s,a+r(s)).
\]

\begin{proposition}[Réversibilité du couplage]
\label{v:rec:completaciones:statement:05}
L’application \(U_r\) est une permutation et
\[
 U_r^{-1}(s,a)=(s,a-r(s)).
\]
La perte d’information n’apparaît qu’après l’application d’un lecteur au
registre ou le conditionnement de l’état par le résultat.
\end{proposition}

\begin{proof}
La formule proposée pour l’inverse satisfait les deux compositions par la
loi de groupe de \(\mathbb Z/9\mathbb Z\).
\end{proof}

Cette proposition formalise la définition opératoire de la mesure : coupler,
lire et restreindre les futurs compatibles sont trois applications différentes. Le
couplage complet est réversible ; le quotient observationnel peut ne pas l’être.

\subsection{Correspondances algébriques et réalisations physiques}
\label{v:rec:completaciones:section:11}

Les complétions de Fock, l’exclusion conditionnée à la parité, le secteur
harmonique de APP, la loi de Gauss discrète et la réversibilité de l’appareil sont des théorèmes
mathématiques complets. Leurs identifications physiques forment un diagramme qui
doit commuter :
\[
 \text{parité de feuillet}
 \longrightarrow\text{échange gradué}
 \longrightarrow\text{CAR/CCR}
 \longrightarrow\text{statistique physique},
\]
\[
 H^1(\mathcal A_9)\simeq\ker\Delta_1
 \hookrightarrow C^1/\operatorname{im}d_0
 \longrightarrow\text{réalisation photonique}.
\]
Dans la première ligne, les deux premières flèches possèdent des constructions
explicites ; la dernière requiert la dynamique et les conditions d’équilibre
correspondantes. Dans la seconde, l’identification de Hodge et l’inclusion
des classes harmoniques dans l’espace des connexions modulo jauge sont les
constructions de la section~\ref{v:rec:completaciones:section:08}. La valeur propre
zéro est celle du laplacien cellulaire : identifier ces modes à des excitations
sans masse et aux deux polarisations locales du photon requiert en outre la
dynamique, la relation de dispersion et la réalisation physique de jauge qui y sont
indiquées. Cette séparation permet de développer les ponts sans confondre une
analogie de cardinalité avec une représentation.
